\subsection{Conjuntos}

Seja $A$ um conjunto. O conjunto potência de $A$, ou conjunto das partes,
é denotado por $2^{A}$. Seja $S$ um conjunto de subconjuntos de $2^{A}$ 
($S \subseteq 2^{A}$); escrevemos:

$$\bigcup S = \bigcup_{X \in S} X$$

\subsection{Relações}

Seja $a, b \in A$ e seja $\ltimes$ uma relação binária em $A$, ou seja,
$\ltimes \subseteq A \times A$. Usamos a notação $a \ltimes b$ para expressar
o elemento $(a, b) \in \ltimes$.

Denotamos por $\rtimes$ a relação inversa associada a $\ltimes$, isto é,
$(a \rtimes b) \iff (b \ltimes a)$.

Denotamos o conjunto de elementos tais que a relação
$a \ltimes b$ é verdadeira por:

$$\ltimes(a) = \{ b \in A : a \ltimes b\}$$

\subsection{Funções}

Seja $B$ um conjunto. Uma função (ou uma aplicação) $\mathcal{F}$ de $A$
para $B$ é uma parte de $A \times B$ tal que, para qualquer $a \in A$, existe
exatamente um $b \in B$ tal que $(a, b) \in \mathcal{F}$. Escrevemos
$b = \mathcal{F}(a)$ para expressar tal par ordenado. Denotamos por
$B^{A}$ o conjunto de todas as funções que levam $A$ para $B$.

Seja $\widetilde{A} \subseteq A$. Define-se o conjunto
$\mathcal{F}(\widetilde{A}) = \{\mathcal{F}(a) : a \in \widetilde{A}\}$.
Quando $\mathcal{F}(\widetilde{A}) = \{b\}$ escrevemos
$\mathcal{F}(\widetilde{A}) = b$ por abuso de notação. Seja
$\widetilde{B} \subseteq B$, escrevemos
$\mathcal{F}^{-1}(\widetilde{B}) = \{a \in A : \mathcal{F}(a) \in \widetilde{B}\}$.
Quando $\widetilde{B} = \{b\}$ escrevemos $\mathcal{F}^{-1}(b)$ ao invés de
$\mathcal{F}^{-1}(\{b\})$ por abuso de notação.

\subsection{Ordem}
Seja $\preccurlyeq$ uma relação binária. Dizemos que $\preccurlyeq$ é uma ordem
em $A$, mais precisamente que $(A, \preccurlyeq)$ é um conjunto ordenado, se
a relação é reflexiva, transitiva e antissimétrica. Dizemos que $\preccurlyeq$
é totalmente ordenado se, além de ser ordenado, também possui a condição de
comparabilidade (ou totalidade). Dizemos que $\preccurlyeq$ é uma ordem
parcial em $A$, ou \textit{poset}, se não possui a condição de comparabilidade.

Suponha que $(A, \preccurlyeq)$ é um conjunto ordenado. Seja
$\widetilde{A} \subseteq A$. Define-se o seguinte:

\begin{itemize}
    \item O supremo de $(\widetilde{A}, \preccurlyeq)$ é denotado por $\supremo{\preccurlyeq}\widetilde{A}$.
    \item O ínfimo de $(\widetilde{A}, \preccurlyeq)$ é denotado por $\infimo{\preccurlyeq}\widetilde{A}$.
    \item O conjunto dos elementos maximais de $(A, \preccurlyeq)$ é denotado por $\maximal{\preccurlyeq}\widetilde{A}$.
    \item O conjunto dos elementos minimais de $(A, \preccurlyeq)$ é denotado por $\minimal{\preccurlyeq}\widetilde{A}$.

\end{itemize}

O diagrama de Hasse do conjunto ordenado $(A, \preccurlyeq)$ é denotado por
$(A, \lhd)$ e corresponde à remoção de redundâncias das propriedades reflexivas
e transitivas. Por abuso de notação, o diagrama de Hasse no conjunto ordenado
$(\widetilde{A}, \preccurlyeq)$ também é denotado por $(\widetilde{A}, \lhd)$.

Quando $(A, \preccurlyeq)$ é um conjunto totalmente ordenado e $a \preccurlyeq b$
com $a, b \in A$, denotamos o intervalo de elementos entre $a$ e $b$ por

$$\llbracket a, b \rrbracket = \{c \in A : a \preccurlyeq c \preccurlyeq b\}$$

\subsection{Árvores}
\label{sec:arvores}
Seja $A$ um conjunto não vazio e $\preccurlyeq$ uma ordem em $A$. O conjunto
ordenado $(A, \preccurlyeq)$ é considerado de ordem hierárquica superior
se as seguintes condições são satisfeitas:

\begin{enumerate}
    \item $(A, \preccurlyeq)$ admite máximo $\supremo{\preccurlyeq}\widetilde{A}$.
    \item $\forall a \in A$, o conjunto $(\preccurlyeq(a), \preccurlyeq)$ é totalmente ordenado.
\end{enumerate}

Sob essas condições, podemos dizer que, se $(A, \preccurlyeq)$ é uma ordem hierárquica
superior, o diagrama de Hasse $(A, \lhd)$ de tal conjunto ordenado é uma árvore em que:

\begin{enumerate}
    \item A raiz é $\supremo{\preccurlyeq} A$.
    \item O conjunto de folhas é $\minimal{\preccurlyeq} A$.
    \item O conjunto de arestas é $\lhd$.
\end{enumerate}

A relação $\lhd$ pode ser vista como uma função $\lhd : A \rightarrow A$ e,
portanto, podemos escrever $\lhd(a) = b$ ao invés de $a \lhd b$.

Seja $\emptyset \subset \widetilde{A} \subseteq A$. Uma subárvore de $(A, \lhd)$ é definida por:

$$\overset{\triangle}{A} = \bigcup_{a \in \widetilde{A}} \llbracket a,  \supremo{\preccurlyeq} \widetilde{A} \rrbracket$$

Ou seja, a árvore $(\overset{\triangle}{A}, \lhd)$ é uma subárvore de
$(A, \lhd)$. Uma propriedade importante é: se $\overset{\triangle}{A} \ne A$
então:

\begin{itemize}
    \item $(A \setminus \overset{\triangle}{A}, \preccurlyeq)$ continua sendo um conjunto ordenado superior.
    \item $(A \setminus \overset{\triangle}{A}, \lhd)$ continua sendo uma árvore.
\end{itemize}

\subsection{Árvore de componentes}

Com as definições de conjuntos, relações, funções, ordem e árvores agora podemos
definir o que é uma imagem e o que é a árvore de componentes dessa imagem.

\begin{definition}[relação de adjacência]
    Seja $\mathbb{P}$ um conjunto finito. Uma relação binária {\Large $\frown$}$\subseteq \mathbb{P} \times \mathbb{P}$
    é dita de adjacência se for irreflexiva e simétrica.
\end{definition}

\begin{definition}[relação de conectividade]
    Uma relação de adjacência {\Large $\frown$} induz uma relação de conectividade $\sim$ em
    $\mathbb{P}$. Dois elementos $x, y \in \mathbb{P}$ são ditos conexos ou conectados ($x \sim y$)
    se e somente se $x \frown y$ ou se existe uma sequência finita com $k \ge 1$
    formada por elementos $x_0, x_1, \dots, x_k \in \mathbb{P}$ tal que:

    \begin{enumerate}
        \item $x_0 = x$.
        \item $x_k = y$.
        \item $x_i \frown x_{i+1}$ para todo $i \in \{0, 1, \dots, k-1\}$.
    \end{enumerate}
\end{definition}

\begin{definition}
    Um conjunto $X \subseteq \mathbb{P}$ é dito ser uma componente conexa
    se for maximal relativamente à propriedade de conexidade.
\end{definition}

Para qualquer $X \subseteq \mathbb{P}$, o conjunto de componentes conexas de $X$
é denotado por $\Pi(X)$ (com $\Pi(\emptyset)=\emptyset$). Note que,
quando $X \ne \emptyset$, o conjunto $\Pi(X)$ é uma partição de $X$.
Quando $X \ne \emptyset$, usamos $\mathscr{C}(X)$ para indicar que $X$ é conexo, isto é, que possui apenas uma componente conexa
($\Pi(X)=\{X\}$).

\begin{definition}
Seja $\mathbb{V} \subseteq \mathbb{Z}$ um conjunto finito com a relação de ordem
$\le$. Seja $\bot$ e $\top$ os menores e maiores elementos de $\mathbb{V}$, respectivamente.
Uma imagem em níveis de cinza $I$ é uma função $I:\mathbb{P}\rightarrow\mathbb{V}$, onde
o domínio $\mathbb{P}$ (comumente chamado de \textit{pixels}) possui relação de adjacência e o contradomínio
corresponde aos níveis de cinza da imagem.
\end{definition}

Seja $v \in \mathbb{V}$. A função de limiar para um valor $v$ é definida como

\begin{equation}
    \left| \begin{aligned}
        \lambda_v : \quad \mathbb{V}^{\mathbb{P}} & \longrightarrow 2^{\mathbb{P}} \\
        \mathcal{F} & \longrightarrow \{x \in \mathbb{P} : v \le \mathcal{F}(x)\}
    \end{aligned} \right.
    \label{eq:limiar}
\end{equation}

O limiar para um nível de cinza $v$ para uma imagem $I$ é definido como

$$\Lambda_v = \lambda_v(I)$$

O conjunto de componentes conexas é

$$\Theta_v = \Pi (\Lambda_v)$$

Seja o intervalo $\llbracket \alpha, \beta \rrbracket \subseteq \mathbb{V}$,
o conjunto de todas as componentes conexas para os valores de $v$ é definido por

$$\Theta_{\llbracket \alpha, \beta \rrbracket} = \bigcup_{v \in \llbracket \alpha, \beta \rrbracket} \Theta_v$$

Portanto, o conjunto de todas as componentes conexas para todos os limiares
é denotado por 

$$\Theta = \Theta_{\llbracket \bot, \top \rrbracket}$$

A partir da definição de $\Theta$ agora podemos introduzir o conceito de
árvore de componentes. Considere a relação binária parcial $\subseteq$ em 
$\Theta$, ou seja, considere o conjunto de ordem parcial $(\Theta, \subseteq)$.
Note que tal conjunto é considerado de ordem hierárquica superior
devido ao conceito apresentado em \ref{sec:arvores}, pois:

\begin{enumerate}
    \item O máximo de tal conjunto corresponde a $\mathbb{P} = \supremo{\subseteq} \Theta$.
    \item Para qualquer componente $X_v \in \Theta_v$ ($v \in \llbracket \bot + 1, \top \rrbracket$)
    existe exatamente uma componente $X_{v-1} \in \Theta_{v-1}$ tal que $X_v \subseteq X_{v-1}$.
    Portanto, $(\subseteq(X), \subseteq)$ é totalmente ordenado.
\end{enumerate}

Conforme apresentado em \ref{sec:arvores}, o diagrama de Hasse de 
$(\Theta, \subseteq)$ é uma árvore. Como essa árvore construída armazena a
informação das componentes conexas, a ela denominamos de árvore de componentes;
mais especificamente, a árvore construída é a \textit{max tree}.

\begin{definition}
    A \textit{max tree} $\mathscr{T}$ é o diagrama de Hasse $(\Theta, \lhd)$ do conjunto
    de ordem parcial $(\Theta, \subseteq)$. Os elementos de $\Theta$ são
    os nós da árvore. A raiz é $\supremo{\subseteq}\Theta$. As folhas são
    os elementos de $\minimal{\subseteq}\Theta$. As arestas são os
    elementos de $\lhd$.
\end{definition}

A título de curiosidade, a \textit{max tree} possui uma árvore dual chamada
\textit{min tree}; para sua construção basta substituir o operador
$\le$ por $\ge$ na equação \eqref{eq:limiar}.

É fácil notar que essa definição de \textit{max tree} não é computacionalmente eficiente em termos de memória.
Em \cite{luz_alves_component_2026} é adotada uma versão mais eficiente. Como
o propósito deste projeto de pesquisa é apenas contextualizar o problema, tal assunto não será abordado aqui.

\subsection{Merge}
Conforme explicado em \ref{sec:introducao}, a árvore de componentes é amplamente
utilizada para diversas operações. No entanto, quando a imagem é alterada
faz-se necessária a reconstrução total da árvore. Nesse sentido, \cite{luz_alves_component_2026}
propõe a operação de mescla (\textit{merge}) entre os nós da árvore. A seguir
formular-se-á o problema da mescla com o intuito de fornecer o contexto para o objetivo desta pesquisa.

\textbf{Pergunta 1}: sejam $I^{-}: \mathbb{P} \rightarrow \mathbb{V}$ uma imagem antes de sofrer
alguma transformação, $\mathscr{T}^{-}_{\text{max}}$ a \textit{max tree} associada a $I^{-}$,
$C \subset \mathbb{P}$ um conjunto conexo, $a = \infimo{\le}I^{-}(C)$ e
$b \in \mathbb{V}$ tal que $b > \supremo{\le}I^{-}(C)$.
A imagem $I^{+}$ é definida como:

$$\begin{array}{|llll}
   I^{+} : & \mathbb{P} & \rightarrow & \mathbb{V}  \\
           & x \in \mathbb{P} \setminus C & \rightarrow & I^{-}(x) \\
           & x \in C & \rightarrow & b
\end{array}$$

Seja $\mathscr{T}_{\text{max}}^{+}$ a \textit{max tree} associada a $I^{+}$.
Sob essas condições, como devemos modificar $\mathscr{T}_{\text{max}}^{-}$
para obter $\mathscr{T}_{\text{max}}^{+}$?

\textbf{Pergunta 2}: sejam $\mathscr{T}_{\text{min}}^{-}$ a \textit{min tree} de $I^{-}$
 e $\overset{\triangle}{\mathscr{T}_{\text{min}}^{-}}$ uma subárvore
 dessa \textit{min tree}. Seja $\mathscr{T}_{\text{min}}^{+}$
 a \textit{min tree} obtida ao remover a subárvore $\overset{\triangle}{\mathscr{T}_{\text{min}}^{-}}$
 de $\mathscr{T}_{\text{min}}^{-}$. Seja $I^{+}$ a imagem reconstruída
 a partir de $\mathscr{T}_{\text{min}}^{+}$. Como modificar
 $\mathscr{T}_{\text{max}}^{-}$ para obter $\mathscr{T}_{\text{max}}^{+}$ sob
 essas condições?
 
\begin{figure}[ht]
    \centering
    % Subfigura (a)
    \begin{subfigure}[b]{0.48\textwidth}
        \centering
        \begin{tikzpicture}
            \begin{axis}[
                width=\textwidth,
                height=8cm,
                ymin=0, ymax=9.5,
                xmin=0.5, xmax=15.5,
                xtick={1,2,3,4,5,6,7,8,9,10,11,12,13,14,15},
                xticklabels={$\mathbf{x}_1$, $\mathbf{x}_2$, $\mathbf{x}_3$, $\mathbf{x}_4$, $\mathbf{x}_5$, $\mathbf{x}_6$, $\mathbf{x}_7$, $\mathbf{x}_8$, $\mathbf{x}_9$, $\mathbf{x}_{10}$, $\mathbf{x}_{11}$, $\mathbf{x}_{12}$, $\mathbf{x}_{13}$, $\mathbf{x}_{14}$, $\mathbf{x}_{15}$},
                ytick={0,1,2,3,4,5,6,7,8,9},
                ymajorgrids=true,
                grid style={dotted, gray!80},
                axis x line=bottom,
                axis y line=left,
                enlarge x limits=false,
                % Posicionando os labels dos eixos como na imagem
                every axis y label/.style={at={(0,1)}, anchor=south, yshift=2pt},
                ylabel={$\mathbb{V}$},
                every axis x label/.style={at={(1,0)}, anchor=west, xshift=2pt},
                xlabel={$\mathbb{P}$},
                % Largura das barras para que fiquem coladas
                bar width=14.5pt,
            ]
            
            % Barras pretas
            \addplot[ybar, fill=black, draw=black] coordinates {
                (1,0.2) (2,7) (3,4) (4,8) (5,9) (6,5) (7,4) (8,2) (9,5) (10,3) (11,8) (12,7) (13,6) (14,1) (15,4)
            };
            
            % Marcadores pretos (fora do conjunto C)
            \addplot[only marks, mark=square*, mark options={draw=white, fill=black, scale=1.3, line width=0.5pt}] coordinates {
                (1,0.2) (2,7) (3,4) (4,8) (5,9) (11,8) (12,7) (13,6) (14,1) (15,4)
            };
            
            % Marcadores rosas (dentro do conjunto C)
            \addplot[only marks, mark=square*, mark options={draw=white, fill=mypink, scale=1.3, line width=0.5pt}] coordinates {
                (6,5) (7,4) (8,2) (9,5) (10,3)
            };
            \end{axis}
        \end{tikzpicture}
        \caption{$I^-$}
    \end{subfigure}
    \hfill
    % Subfigura (b)
    \begin{subfigure}[b]{0.48\textwidth}
        \centering
        \begin{tikzpicture}
            \begin{axis}[
                width=\textwidth,
                height=8cm,
                ymin=0, ymax=9.5,
                xmin=0.5, xmax=15.5,
                xtick={1,2,3,4,5,6,7,8,9,10,11,12,13,14,15},
                xticklabels={$\mathbf{x}_1$, $\mathbf{x}_2$, $\mathbf{x}_3$, $\mathbf{x}_4$, $\mathbf{x}_5$, $\mathbf{x}_6$, $\mathbf{x}_7$, $\mathbf{x}_8$, $\mathbf{x}_9$, $\mathbf{x}_{10}$, $\mathbf{x}_{11}$, $\mathbf{x}_{12}$, $\mathbf{x}_{13}$, $\mathbf{x}_{14}$, $\mathbf{x}_{15}$},
                ytick={0,1,2,3,4,5,6,7,8,9},
                ymajorgrids=true,
                grid style={dotted, gray!80},
                axis x line=bottom,
                axis y line=left,
                enlarge x limits=false,
                every axis y label/.style={at={(0,1)}, anchor=south, yshift=2pt},
                ylabel={$\mathbb{V}$},
                every axis x label/.style={at={(1,0)}, anchor=west, xshift=2pt},
                xlabel={$\mathbb{P}$},
                bar width=14.5pt,
            ]
            
            % Barras pretas
            \addplot[ybar, fill=black, draw=black] coordinates {
                (1,0.2) (2,7) (3,4) (4,8) (5,9) (6,7) (7,7) (8,7) (9,7) (10,7) (11,8) (12,7) (13,6) (14,1) (15,4)
            };
            
            % Marcadores pretos (fora do conjunto C)
            \addplot[only marks, mark=square*, mark options={draw=white, fill=black, scale=1.3, line width=0.5pt}] coordinates {
                (1,0.2) (2,7) (3,4) (4,8) (5,9) (11,8) (12,7) (13,6) (14,1) (15,4)
            };
            
            % Marcadores rosas (dentro do conjunto C ajustado para b=7)
            \addplot[only marks, mark=square*, mark options={draw=white, fill=mypink, scale=1.3, line width=0.5pt}] coordinates {
                (6,7) (7,7) (8,7) (9,7) (10,7)
            };
            \end{axis}
        \end{tikzpicture}
        \caption{$I^+$}
    \end{subfigure}
    \caption{As imagens acima ilustram o cenário da \textbf{pergunta 1}. Exemplo
    retirado de \cite{luz_alves_component_2026}}
\end{figure}

\begin{figure}[ht]
    \centering
    % Subfigura (a)
    \begin{subfigure}[b]{0.48\textwidth}
        \centering
        \begin{tikzpicture}
            \begin{axis}[
                width=\textwidth,
                height=8cm,
                ymin=0, ymax=9.5,
                xmin=0.5, xmax=15.5,
                xtick={1,2,3,4,5,6,7,8,9,10,11,12,13,14,15},
                xticklabels={$\mathbf{x}_1$, $\mathbf{x}_2$, $\mathbf{x}_3$, $\mathbf{x}_4$, $\mathbf{x}_5$, $\mathbf{x}_6$, $\mathbf{x}_7$, $\mathbf{x}_8$, $\mathbf{x}_9$, $\mathbf{x}_{10}$, $\mathbf{x}_{11}$, $\mathbf{x}_{12}$, $\mathbf{x}_{13}$, $\mathbf{x}_{14}$, $\mathbf{x}_{15}$},
                ytick={0,1,2,3,4,5,6,7,8,9},
                ymajorgrids=true,
                grid style={dotted, gray!80},
                axis x line=bottom,
                axis y line=left,
                enlarge x limits=false,
                % Posicionando os labels dos eixos como na imagem
                every axis y label/.style={at={(0,1)}, anchor=south, yshift=2pt},
                ylabel={$\mathbb{V}$},
                every axis x label/.style={at={(1,0)}, anchor=west, xshift=2pt},
                xlabel={$\mathbb{P}$}
            ]
            
            % Linhas horizontais cinzas (componentes)
            \draw[gray, thick] (axis cs:1,0) -- (axis cs:15,0);
            \draw[gray, thick] (axis cs:2,1) -- (axis cs:15,1);
            \draw[gray, thick] (axis cs:2,2) -- (axis cs:13,2);
            \draw[gray, thick] (axis cs:3,3) -- (axis cs:7,3);
            \draw[gray, thick] (axis cs:9,3) -- (axis cs:13,3);
            \draw[gray, thick] (axis cs:3,4) -- (axis cs:7,4);
            \draw[gray, thick] (axis cs:9,4) -- (axis cs:13,4);
            \draw[gray, thick] (axis cs:4,5) -- (axis cs:6,5);
            \draw[gray, thick] (axis cs:11,5) -- (axis cs:13,5);
            \draw[gray, thick] (axis cs:4,6) -- (axis cs:5,6);
            \draw[gray, thick] (axis cs:11,6) -- (axis cs:13,6);
            \draw[gray, thick] (axis cs:4,7) -- (axis cs:5,7);
            \draw[gray, thick] (axis cs:11,7) -- (axis cs:12,7);
            \draw[gray, thick] (axis cs:4,8) -- (axis cs:5,8);

            % Arestas da árvore (linhas pretas)
            \draw[thick] (axis cs:8,0) -- (axis cs:9,1);
            \draw[thick] (axis cs:9,1) -- (axis cs:8,2);
            \draw[thick] (axis cs:9,1) -- (axis cs:15,4);
            \draw[thick] (axis cs:8,2) -- (axis cs:5,4);
            \draw[thick] (axis cs:8,2) -- (axis cs:11,3);
            \draw[thick] (axis cs:5,4) -- (axis cs:2,7);
            \draw[thick] (axis cs:5,4) -- (axis cs:5,5);
            \draw[thick] (axis cs:5,5) -- (axis cs:5,8);
            \draw[thick] (axis cs:5,8) -- (axis cs:6,9);
            \draw[thick] (axis cs:11,3) -- (axis cs:10,5);
            \draw[thick] (axis cs:11,3) -- (axis cs:12,6);
            \draw[thick] (axis cs:12,6) -- (axis cs:11,7);
            \draw[thick] (axis cs:11,7) -- (axis cs:11,8);

            % Nós (marcadores circulares coloridos)
            \addplot[only marks, mark=*, mark options={draw=black, fill=myyellow, scale=1.5}] coordinates {
                (8,0) (9,1) (15,4) (2,7) (6,9) (11,8)
            };
            \addplot[only marks, mark=*, mark options={draw=black, fill=mycyan, scale=1.5}] coordinates {
                (8,2)
            };
            \addplot[only marks, mark=*, mark options={draw=black, fill=myblue, scale=1.5}] coordinates {
                (5,4) (5,5) (10,5) (11,3) (12,6) (11,7)
            };
            \addplot[only marks, mark=*, mark options={draw=black, fill=mygreen, scale=1.5}] coordinates {
                (5,8)
            };
            \end{axis}
        \end{tikzpicture}
        \caption{$\mathscr{T}^-_{\max}$}
    \end{subfigure}
    \hfill
    % Subfigura (b)
    \begin{subfigure}[b]{0.48\textwidth}
        \centering
        \begin{tikzpicture}
            \begin{axis}[
                width=\textwidth,
                height=8cm,
                ymin=0, ymax=9.5,
                xmin=0.5, xmax=15.5,
                xtick={1,2,3,4,5,6,7,8,9,10,11,12,13,14,15},
                xticklabels={$\mathbf{x}_1$, $\mathbf{x}_2$, $\mathbf{x}_3$, $\mathbf{x}_4$, $\mathbf{x}_5$, $\mathbf{x}_6$, $\mathbf{x}_7$, $\mathbf{x}_8$, $\mathbf{x}_9$, $\mathbf{x}_{10}$, $\mathbf{x}_{11}$, $\mathbf{x}_{12}$, $\mathbf{x}_{13}$, $\mathbf{x}_{14}$, $\mathbf{x}_{15}$},
                ytick={0,1,2,3,4,5,6,7,8,9},
                ymajorgrids=true,
                grid style={dotted, gray!80},
                axis x line=bottom,
                axis y line=left,
                enlarge x limits=false,
                every axis y label/.style={at={(0,1)}, anchor=south, yshift=2pt},
                ylabel={$\mathbb{V}$},
                every axis x label/.style={at={(1,0)}, anchor=west, xshift=2pt},
                xlabel={$\mathbb{P}$}
            ]
            
            % Linhas horizontais cinzas (componentes)
            \draw[gray, thick] (axis cs:1,0) -- (axis cs:15,0);
            \draw[gray, thick] (axis cs:2,1) -- (axis cs:15,1);
            \draw[gray, thick] (axis cs:2,2) -- (axis cs:13,2);
            \draw[gray, thick] (axis cs:2,3) -- (axis cs:13,3);
            \draw[gray, thick] (axis cs:2,4) -- (axis cs:13,4);
            \draw[gray, thick] (axis cs:4,5) -- (axis cs:13,5);
            \draw[gray, thick] (axis cs:4,6) -- (axis cs:13,6);
            \draw[gray, thick] (axis cs:4,7) -- (axis cs:12,7);
            \draw[gray, thick] (axis cs:4,8) -- (axis cs:5,8);
            
            % Arestas da árvore (linhas pretas)
            \draw[thick] (axis cs:8,0) -- (axis cs:9,1);
            \draw[thick] (axis cs:9,1) -- (axis cs:15,4);
            \draw[thick] (axis cs:9,1) -- (axis cs:8,4);
            \draw[thick] (axis cs:8,4) -- (axis cs:2,7);
            \draw[thick] (axis cs:8,4) -- (axis cs:9,6);
            \draw[thick] (axis cs:9,6) -- (axis cs:5,8);
            \draw[thick] (axis cs:9,6) -- (axis cs:11,8);
            \draw[thick] (axis cs:5,8) -- (axis cs:6,9);

            % Nós (marcadores circulares coloridos)
            \addplot[only marks, mark=*, mark options={draw=black, fill=myyellow, scale=1.5}] coordinates {
                (8,0) (9,1) (15,4) (2,7) (11,8) (6,9)
            };
            \addplot[only marks, mark=*, mark options={draw=black, fill=mycyan, scale=1.5}] coordinates {
                (8,4)
            };
            \addplot[only marks, mark=*, mark options={draw=black, fill=myblue, scale=1.5}] coordinates {
                (9,6)
            };
            \addplot[only marks, mark=*, mark options={draw=black, fill=mygreen, scale=1.5}] coordinates {
                (5,8)
            };
            \end{axis}
        \end{tikzpicture}
        \caption{$\mathscr{T}^+_{\max}$}
    \end{subfigure}
    \caption{\textit{Max tree} das imagens $I^-$ e $I^+$ respectivamente.
    Exemplo retirado de \cite{luz_alves_component_2026}}
\end{figure}